% System Design Document — build: docs/latex/compile_system_design.sh → docs/reports/System_Design.pdf
\documentclass[11pt,a4paper]{article}

\usepackage[utf8]{inputenc}
\usepackage[T1]{fontenc}
\usepackage{lmodern}
\usepackage[a4paper,margin=1in]{geometry}
\usepackage{microtype}
\usepackage{parskip}
\usepackage[hidelinks]{hyperref}

\title{Exoskeleton Control System\\[0.35em]\large System Design Document}
\author{AutomationForExoskeleton}
\date{Version 1.3 --- April 5, 2026}

\begin{document}
\maketitle

\section{Executive Summary}
This document describes the architecture of the software-centric control module for the ExoTechHK lower-limb exoskeleton. The system integrates inertial measurement units (IMUs), Support Vector Machine (SVM) classification, and optional Kalman Filter visualization to provide real-time, adaptive assistance.

Version 1.3 uses \textbf{HuGaDB} from \textbf{GitHub v1} only, but now treats it as a protocol-aware and corruption-aware dataset: \path{LoadHuGaDB.m} reads raw \path{.txt} files, applies the official gyro corruption matrix (CSV next to the loader), preserves \texttt{single\_activity} versus \texttt{multi\_activity\_sequence} metadata, and stores quality flags in the cache. Binary and multiclass HuGaDB training/evaluation default to excluding \path{ExoConfig.HUGADB.HELDOUT_SUBJECTS}, filtering to \path{ExoConfig.HUGADB.DEFAULT_PROTOCOLS}, and skipping invalid sessions/windows so the policy matches replay simulation. The common feature vector remains \textbf{30-dimensional} (six IMU slots $\times$ five statistics per slot), with \textbf{binary} and \textbf{multiclass} ECOC SVM paths and the asymmetric finite state machine for assist gating.

\section{System Architecture}

\subsection{Configuration (\texttt{ExoConfig.m})}
The system is driven by a configuration class that centralizes tuning parameters.
\begin{itemize}
  \item \textbf{Sampling Rate:} 100 Hz
  \item \textbf{Window Size:} 1.0 s (100 samples)
  \item \textbf{Step Size:} 0.5 s (50 samples)
  \item \textbf{Paths:} \path{usc_had_dataset.mat}, \path{hugadb_dataset.mat}, \path{Binary_SVM_Model_hugadb_only.mat}, \path{Multiclass_SVM_ECOC_usc_had.mat}, \path{Multiclass_SVM_ECOC_hugadb.mat}
  \item \textbf{LOCOMOTION:} \texttt{N\_IMU\_SLOTS = 6}, \texttt{FEATURES\_PER\_IMU = 5} (must match \path{Features.m})
\end{itemize}

\subsection{Data Flow Pipeline}
\begin{enumerate}
  \item \textbf{Input:} Public \path{.mat} caches via loaders; held-out HuGaDB sessions replay through \path{LoadHuGaDBSimulationData.m}.
  \item \textbf{Pre-processing:} HuGaDB cache is normalized to m/s$^2$ and rad/s during loading, tagged by session protocol, and screened by a HuGaDB quality gate before windows are used for training/evaluation.
  \item \textbf{Feature extraction:} Five features per IMU (\path{Features.m}). HuGaDB uses \textbf{six IMUs} per window (\path{FeaturesFromImuStack}). USC-HAD uses \textbf{one IMU + zero padding} (\path{LocomotionFeatureVector}) for benchmark comparisons.
  \item \textbf{Inference --- binary:} RBF SVM (\texttt{fitcsvm}) $\rightarrow$ class 0/1 (inactive vs active movement per dataset label maps).
  \item \textbf{Inference --- multiclass (optional):} ECOC SVM (\texttt{fitcecoc}) on \textbf{native} labels: USC-HAD \textbf{12} trial activities or HuGaDB \textbf{12} per-sample IDs (\path{ActivityClassRegistry.m}). Exo assist maps active-movement classes through \path{RealtimeFsmFromActivityClass.m} (dataset argument: \texttt{'usc\_had'} / \texttt{'hugadb'}).
  \item \textbf{State machine:} \path{RealtimeFsm.m} applies asymmetric hysteresis on \textbf{binary} predictions.
  \item \textbf{Fusion:} Optional Kalman (\texttt{imufilter}) computes a segment-pitch visualization trace in parallel.
  \item \textbf{Output:} Control command (0/1) + visualization trace; multiclass scripts add activity labels for logging/plots.
\end{enumerate}

\section{Core Components}

\subsection{Feature extraction (\texttt{Features.m}, \texttt{FeaturesFromImuStack.m}, \texttt{LocomotionFeatureVector.m})}
Per IMU window, the feature vector is \textbf{1$\times$5}:
\begin{enumerate}
  \item Mean acc magnitude
  \item Variance of acc magnitude
  \item Mean gyro magnitude
  \item Variance of gyro magnitude
  \item Dominant FFT frequency of vertical acceleration
\end{enumerate}

\textbf{Training on public data:} HuGaDB trials store \textbf{N$\times$3$\times$6} acc/gyro (\path{LoadHuGaDB.m}); stacked features are \textbf{1$\times$30}. USC-HAD trials use one IMU $\rightarrow$ \textbf{1$\times$5} inside \path{LocomotionFeatureVector} plus \textbf{25 zeros}.

\subsection{Activity labels (\texttt{ActivityClassRegistry.m})}
Multiclass ECOC uses \textbf{dataset-native} IDs (no cross-dataset unification): USC-HAD \textbf{1\ldots12} per trial, HuGaDB \textbf{1\ldots12} per sample (window label = mode over samples). Binary training uses dataset-specific active / non-active ID groups from \path{ExoConfig.DS}.

\subsection{Sensor fusion (\texttt{FusionKalman.m})}
Static class wrapping MATLAB \texttt{imufilter}. Under the HuGaDB replay path it produces an orientation-derived segment-pitch trace for visualization.

\subsection{Finite State Machine (\texttt{RealtimeFsm.m})}
Asymmetric hysteresis: \textbf{3} consecutive active binary predictions to enter assist; \textbf{5} consecutive inactive predictions to exit.

\section{Execution \& File Organization}

\subsection{Path management}
Main scripts resolve \texttt{projectRoot} locally and add \path{config/}, the active \path{scripts/} folder, and \path{src/} (recursive) as needed. Dataset loader functions remain under \path{data/USC-HAD/}, \path{data/HuGaDB/}, etc., and can be added to the MATLAB path when rebuilding local caches.

\textbf{USC-HAD:} \path{LoadUSCHAD.m} $\rightarrow$ \path{data/USC-HAD/usc_had_dataset.mat} (recursive raw \path{.mat} under \path{USC-HAD_raw/}).

\textbf{HuGaDB:} \path{LoadHuGaDB.m} $\rightarrow$ \path{data/HuGaDB/hugadb_dataset.mat} (GitHub v1 under \path{v1_cleanup_github/} flat layout or legacy nested \path{v1_raw}/flat paths; official gyro corruption matrix zeroes bad IMU gyros; \texttt{readmatrix} treats \texttt{na} as missing when present; all-six-corrupt sessions are retained as accelerometer-only samples so classes such as bicycling / sitting-in-car are not lost; six IMUs per row; optional manifest adds \texttt{huGaDBProvenance} / \textbf{\texttt{huGaDBSessionProtocol}}; accepted sessions carry quality metadata and downstream prep functions skip invalid sessions/windows with logged reasons).

\textbf{Binary training / evaluation:} \path{TrainSvmBinary.m}, \path{EvaluateSvmConfusion.m}, \path{RunSvmDatasetAblation.m} (USC-HAD and HuGaDB).

\textbf{Multiclass:} \path{PrepareTrainingDataMulticlass.m}, \path{TrainSvmMulticlass.m}, \path{EvaluateMulticlassConfusion.m}, \path{RunTrainEvalMulticlass.m}, \path{RunExoskeletonPipelineMulticlass.m}.

\textbf{CUHK-REXYIM (local captures):} \path{ImportData.m} $\rightarrow$ \path{data/CUHK-REXYIM/<activity>/} with \path{Accelerometer.csv}, \path{Gyroscope.csv}, optional \path{Annotation.csv}.

\subsection{Directory structure (abridged)}
\begin{verbatim}
AutomationForExoskeleton/
  config/                 # ExoConfig.m, ActivityClassRegistry.m
  data/USC-HAD/           # LoadUSCHAD.m, usc_had_dataset.mat
  data/HuGaDB/            # LoadHuGaDB.m, hugadb_dataset.mat
  data/CUHK-REXYIM/       # project phone CSVs per activity (ImportData.m)
  models/                 # Binary_SVM_Model*.mat, Multiclass_SVM_ECOC_*.mat
  results/                # figures/ and metrics/ outputs
  scripts/                # Train/eval/pipeline scripts
  src/
    acquisition/          # PrepareTrainingData*.m
    classification/       # RealtimeFsm.m, RealtimeFsmFromActivityClass.m
    features/             # Features.m, FeaturesFromImuStack.m, ...
    fusion/
\end{verbatim}

\end{document}
